\documentclass[11pt,a4paper]{article}

\usepackage[margin=25mm]{geometry}
\usepackage{amsmath,amssymb,amsthm}
\usepackage{booktabs}
\usepackage{enumitem}
\usepackage{hyperref}

\title{An Exact Lexicographic Min-Cost Flow Solver\\
for the Laboratory Assignment Problem}
\author{Meteo}
\date{2026-06-13}

\newtheorem{theorem}{Theorem}
\newtheorem{lemma}{Lemma}

\begin{document}
\maketitle

\begin{abstract}
This document describes the mathematical model, algorithm, implementation,
and verification strategy of a C program for the laboratory assignment
problem.  The program assigns every student to exactly one laboratory, never
exceeds laboratory capacities, and assigns at least one student to every
laboratory with positive capacity.  The core algorithm is an exact
network-flow solver: the default mode minimizes rank sum, rank squared sum,
worst assigned rank, and fill-rate objectives lexicographically without using a
Big-M scalarization.  The implementation additionally compresses equivalent
capacity slots and uses an indexed heap so that the exact algorithm remains
fast on maximum-size instances.  To reduce manual checking work for users, it
automatically writes only the simple assignment file by default, and can
optionally write numerical metrics and review reports when the
\texttt{--reports} option is specified.  For reproducibility and performance
inspection, the implementation also provides preset rank-cost and weighted
objective files, objective-printing helpers, exact counterfactual
explanations, an optional solver profile TSV, and process-parallel portfolio
execution.
\end{abstract}

\section{Problem Statement}

Let
\[
  S=\{1,\ldots,N\}, \qquad L=\{1,\ldots,M\}
\]
be the set of students and laboratories.  The target benchmark constraints are
\[
  1 \leq N \leq 1024, \qquad 1 \leq M \leq 256.
\]
The implementation does not hard-code these two benchmark limits as rejection
conditions.  Larger feasible inputs are accepted subject to memory, time, and
standard-C integer-index limits.
Laboratory \(j\) has capacity \(c_j \geq 0\), and the total capacity satisfies
\[
  \sum_{j=1}^{M} c_j \geq N.
\]
Student \(i\) lists \(R\) preferred laboratories.  If laboratory \(j\) is the
\(k\)-th preference of student \(i\), define \(r_{ij}=k\).  If \(j\) is not
listed by student \(i\), this implementation uses the outside-preference rank:
\[
  r_{ij}=M+1.
\]

The decision variable is
\[
  x_{ij} \in \{0,1\},
\]
where \(x_{ij}=1\) means that student \(i\) is assigned to laboratory \(j\).
The hard constraints are
\[
  \sum_{j=1}^{M} x_{ij} = 1 \quad (\forall i\in S),
\]
\[
  \ell_j \leq \sum_{i=1}^{N} x_{ij} \leq c_j \quad (\forall j\in L),
\]
where
\[
  \ell_j =
  \begin{cases}
    1 & (c_j>0),\\
    0 & (c_j=0).
  \end{cases}
\]
If the number of positive-capacity laboratories is greater than \(N\), these
constraints are infeasible and the program reports an input error.

\section{Evaluation Metrics}

For a feasible assignment \(x\), define
\[
  Q(x)=\max_{i,j:x_{ij}=1} r_{ij},
\]
\[
  R_1(x)=\sum_{i=1}^{N}\sum_{j=1}^{M} r_{ij}x_{ij},
  \qquad
  R_2(x)=\sum_{i=1}^{N}\sum_{j=1}^{M} r_{ij}^2x_{ij}.
\]
Minimizing \(R_1(x)\) is equivalent to minimizing the average assigned rank.
For assignments with fixed \(R_1(x)\), minimizing \(R_2(x)\) is equivalent to
minimizing the rank variance because
\[
  \sigma^2(x)
  =\frac{R_2(x)}{N}
  -\left(\frac{R_1(x)}{N}\right)^2.
\]

For positive-capacity laboratories, let \(n_j=\sum_i x_{ij}\).  The minimum
fill rate is
\[
  U_{\min}(x)=\min_{j:c_j>0}\frac{n_j}{c_j}.
\]
The average fill rate is
\[
  U_{\mathrm{avg}}(x)
  = \frac{1}{|\{j:c_j>0\}|}\sum_{j:c_j>0}\frac{n_j}{c_j}.
\]

\section{Optimization Order}

The evaluation setup does not define one absolute scalar score.  It gives
five separate metrics whose relative rankings are averaged.  Therefore the
default solver mode uses a rubric-aligned deterministic Pareto-lexicographic
objective:
\[
  \left(
    R_1,\ R_2,\ Q,\ -U_{\mathrm{avg}},\ -U_{\min}
  \right).
\]
All five components are optimized exactly under this order.  This does not
claim to predict an external final average ranking against other solvers,
because other solvers' outputs are unknown.  It does guarantee that the program
does not return a solution that is worse under the stated lexicographic
objective.

The implementation also keeps the earlier conservative fairness order as
\texttt{--objective fair}:
\[
  Q,\ R_1,\ R_2,\ -U_{\min},\ -U_{\mathrm{avg}}.
\]
The \texttt{balanced} mode uses
\[
  R_1,\ R_2,\ Q,\ -U_{\min},\ -U_{\mathrm{avg}},
\]
and \texttt{guarded} first bounds \(Q\) by the minimum feasible value plus a
user-specified slack, then optimizes the default rubric order inside that
guard.

For practical student satisfaction, the implementation also provides a
\texttt{satisfaction} objective mode.  This mode replaces the raw rank \(r\)
by a non-linear dissatisfaction cost \(d(r)\), then optimizes
\[
  D_1,\ D_2,\ Q,\ -U_{\mathrm{avg}},\ -U_{\min},
\]
where
\[
  D_1=\sum_i d(r_i), \qquad D_2=\sum_i d(r_i)^2.
\]
The default model is
\[
  d(1)=0,\quad
  d(r)=100+30(r-2)+5(r-2)^2\ (2\le r\le M),\quad
  d(M+1)=10000.
\]
These constants can be changed from the command line or by a rank-cost table.
Because \(D_1\) and \(D_2\) are additive over student--laboratory edges, this
mode remains a single lexicographic min-cost-flow optimization followed by the
same exact \(Q\), \(U_{\mathrm{avg}}\), and \(U_{\min}\) tie-breaks.  It does
not require the two-dimensional threshold enumeration used by weighted exact
optimization.

The rank-cost table can be supplied from a plain text preset file.  The
repository includes presets such as \texttt{student\_friendly},
\texttt{balanced\_satisfaction}, and \texttt{strict\_outside\_avoidance}.
Weighted exact mode similarly accepts a plain text weight file.  These files
do not change the solver; they only make the exact objective configuration
auditable and easy to reproduce.

The \texttt{fill-convex} mode adds another exact candidate for operational
balance.  Let \(g_j=\min(c_j,N)\) be the graph capacity and
\(G=\sum_j g_j\).  The capacity-proportional ideal count is
\[
  t_j = N\frac{g_j}{G}.
\]
The mode minimizes a scalar objective combining the ordinary rank cost with the
separable convex
penalty
\[
  \sum_j (n_jG-Ng_j)^2.
\]
It is therefore not a rank-first tie-breaker; changing the convex penalty can
change the selected assignment even when the pure rank tuple differs.
This is still a min-cost-flow problem: for the \(x\)-th student assigned to
laboratory \(j\), the marginal penalty is
\[
  G^2(2x-1)-2GNg_j.
\]
The program expands the laboratory-to-sink capacity into one-person edges with
these marginal costs.  Since the marginal-cost representation is exactly the
difference of consecutive squared penalties, the selected assignment is exact
for this convex-balance objective.

\section{Intuitive Explanation of Optimality}

The main reason this method is optimal is that the program does not guess
assignments one by one.  Instead, it changes the assignment problem into a
pipe network.

Think of every student as one unit of water.  The water must leave the source,
pass through exactly one student node, pass into exactly one laboratory node,
and finally reach the sink.  A laboratory pipe has limited capacity, so too
many students cannot enter that laboratory.  A positive-capacity laboratory
also has required slots, so the solver is forced to use at least one slot
there whenever that is mathematically possible.

This means the network has the same choices as the original assignment
problem:
\begin{itemize}[leftmargin=2em]
  \item every valid assignment is one complete integer flow;
  \item every complete integer flow is one valid assignment;
  \item no other hidden choices exist.
\end{itemize}
Therefore, optimizing the flow is the same as optimizing the assignment.

The default mode attaches a small report card to each possible pipe.  The
report card is not a single fragile number; it is an ordered tuple:
\[
  (\text{rank sum},\ \text{rank-square sum},\ \text{worst rank},
  \ \text{average fill rate},\ \text{minimum fill rate}).
\]
The solver compares these tuples exactly from left to right.  Thus it first
minimizes the average rank, then minimizes rank unfairness, then minimizes the
worst assigned rank, and only then uses exact fill-rate tie-breaking.
This is the same as saying:
\begin{quote}
  First make the total rank as small as possible.  If there is still a tie,
  make the spread of ranks as small as possible.  If there is still a tie,
  improve the worst assigned rank, then the average laboratory fill rate, and
  finally the worst laboratory fill rate.
\end{quote}
Because min-cost flow globally minimizes this tuple over all complete flows,
the chosen assignment is not merely locally good.  It is the best assignment
under the stated lexicographic objective.

There is one important fast path in the fair-mode route.  After the worst-rank
threshold has been fixed, if every active student has exactly the same rank
table, then exchanging two students never changes the fair-mode objective
components.  The only information that matters is the count \(n_j\) assigned
to each laboratory.  In that special case, the program skips the
student-laboratory flow and optimizes the count vector directly: it fills
lower-rank laboratories before higher-rank ones, then maximizes the exact
minimum fill rate, and only then maximizes the exact average fill rate.  Other
rank-first modes still remain exact, but they use the general grouped
min-cost-flow path rather than this count-only shortcut.

\section{Formal Correctness Lemmas}

The exactness argument used by the implementation can be audited through the
following lemmas.

\begin{lemma}[Assignment-flow equivalence]
Every feasible assignment corresponds to one integral feasible flow in the
student-laboratory network, and every integral feasible flow corresponds to
one feasible assignment.
\end{lemma}

\begin{proof}
The source sends one unit through each assigned student, each student node has
capacity one from the source, and each unit then enters exactly one laboratory
node before reaching the sink.  Laboratory-to-sink bounds are exactly the
minimum-occupancy and capacity constraints.  Thus the network choices and the
assignment choices are the same.
\end{proof}

\begin{lemma}[Integrality]
The flow solution can be taken to be integral.
\end{lemma}

\begin{proof}
All capacities and lower bounds are integral.  Standard max-flow and
min-cost-flow algorithms on such networks return integral augmentations, so
the resulting flow values are integral.
\end{proof}

\begin{lemma}[Threshold monotonicity]
If rank threshold \(T\) is feasible, every threshold \(T'>T\) is feasible.
\end{lemma}

\begin{proof}
Increasing the threshold only adds student-to-laboratory edges.  It never
removes a feasible flow.
\end{proof}

\begin{lemma}[Binary search returns minimum \(Q\)]
Binary search over threshold feasibility returns the minimum possible worst
assigned rank \(Q\).
\end{lemma}

\begin{proof}
By threshold monotonicity, feasible thresholds form a suffix of
\(\{1,\ldots,M+1\}\).  Binary search returns the first element of this suffix.
\end{proof}

\begin{lemma}[Tuple costs optimize \(R_1\) and \(R_2\)]
With \(Q\) fixed, lexicographic tuple min-cost flow minimizes \(R_1\), and
among those assignments minimizes \(R_2\), without Big-M scalarization.
\end{lemma}

\begin{proof}
The second tuple component is exactly \(R_1\), and the third tuple component
is exactly \(R_2\).  Lexicographic comparison minimizes them in that order.
\end{proof}

\begin{lemma}[Rank variance]
When \(R_1\) is fixed, minimizing \(R_2\) minimizes rank variance.
\end{lemma}

\begin{proof}
The variance is
\[
  \sigma^2=\frac{R_2}{N}-\left(\frac{R_1}{N}\right)^2.
\]
For fixed \(R_1\) and \(N\), only \(R_2\) can change.
\end{proof}

\begin{lemma}[\(U_{\min}\) candidates and monotonicity]
Every achievable minimum fill rate is some \(a/c_j\), and feasibility for a
candidate \(q\) is monotone downward.
\end{lemma}

\begin{proof}
For each positive-capacity laboratory, \(n_j\) is integral, so
\(n_j/c_j=a/c_j\).  If all laboratories satisfy \(n_j/c_j\ge q\), they also
satisfy every smaller candidate.
\end{proof}

\begin{lemma}[Exact \(U_{\mathrm{avg}}\) comparison]
Least-common-denominator integer comparison orders average fill rates exactly.
\end{lemma}

\begin{proof}
Let \(L=\operatorname{lcm}\{c_j\mid c_j>0\}\).  This denominator is fixed for
the input, so comparing
\(\sum_j n_jL/c_j\) is equivalent to comparing
\(\sum_j n_j/c_j\).  The implementation evaluates this integer comparison
with big integers, so no rounding is introduced.
\end{proof}

\begin{lemma}[Equal-capacity bucket compression]
Combining laboratories with the same positive capacity preserves
\(U_{\mathrm{avg}}\) exactly.
\end{lemma}

\begin{proof}
For \(C=\{d\mid \exists j,\ c_j=d,\ d>0\}\),
\[
  \sum_{j:c_j>0}\frac{n_j}{c_j}
  =
  \sum_{d\in C}\frac{\sum_{j:c_j=d}n_j}{d}.
\]
This only combines equal-denominator terms.
\end{proof}

\begin{lemma}[Active rank-row grouping]
After the optimal worst rank \(Q\) is fixed, grouping students with identical
active rank rows preserves all five objective values.
\end{lemma}

\begin{proof}
For a fixed \(Q\), inactive edges with \(r_{ij}>Q\) cannot be used by any
rank-optimal assignment.  If two students have the same available
laboratories and the same rank value on every active edge, swapping their
assigned active laboratories does not change \(Q\), \(R_1\), \(R_2\),
laboratory counts, \(U_{\min}\), or \(U_{\mathrm{avg}}\).  Therefore only
group-to-lab counts matter inside such a group.
\end{proof}

\begin{lemma}[Allowed-set feasibility grouping]
For a fixed threshold \(T\), grouping students with the same allowed
laboratory set preserves feasibility.
\end{lemma}

\begin{proof}
The threshold feasibility test only asks whether each student can use one
laboratory from \(\{j\mid c_j>0,\ r_{ij}\leq T\}\).  Rank values inside this
set are irrelevant until the later min-cost-flow phases.  Therefore students
with the same allowed set are interchangeable for the feasibility network.
A group node with capacity equal to the group size can send exactly the same
number of assignment units to the same set of laboratories as the individual
student nodes could.
\end{proof}

\begin{lemma}[Capacity slot compression]
Replacing identical unit slots by one capacity edge preserves feasible flows
and costs.
\end{lemma}

\begin{proof}
Sending \(k\) units through one capacity-\(c\) edge with linear cost is
equivalent to sending those \(k\) units through \(k\) identical unit edges.
\end{proof}

\section{Flow Model}

The assignment is represented as a bipartite flow network.
\begin{itemize}[leftmargin=2em]
  \item Source \(s\) connects to each student \(i\) with capacity 1.
  \item Student \(i\) connects to laboratory \(j\) with capacity 1 if the
        current rank threshold permits \(r_{ij}\).
  \item Laboratory \(j\) connects to sink \(t\) through capacity slots.
\end{itemize}
Because all capacities are integers, integral min-cost flow returns integral
flows.  Thus every unit of flow corresponds exactly to assigning one student
to one laboratory.

\section{Rank-First Rubric Objective}

The default \texttt{rubric} mode follows the evaluation-metric order:
\[
  R_1 \rightarrow R_2 \rightarrow Q
  \rightarrow U_{\mathrm{avg}} \rightarrow U_{\min}.
\]
It first solves a rank-cost min-cost-flow problem over all rank thresholds up
to \(M+1\), obtaining the minimum pair \((R_1^*,R_2^*)\).  Then it searches the
compressed rank-threshold candidate list for the smallest \(Q\) that can still
achieve \((R_1^*,R_2^*)\).  Finally it fixes the rank pair and \(Q\), maximizes
the exact average fill rate, and among those solutions maximizes the exact
minimum fill rate.

\begin{theorem}
The default \texttt{rubric} solution is lexicographically optimal for
\((R_1,R_2,Q,-U_{\mathrm{avg}},-U_{\min})\).
\end{theorem}

\begin{proof}
The first min-cost-flow run uses student--laboratory edge cost
\((0,r_{ij},r_{ij}^2)\), so its second and third total-cost components are
exactly \(R_1\) and \(R_2\).  Integral min-cost flow therefore returns the
minimum achievable rank pair.  Restricting to thresholds \(Q\) and re-solving
with the rank pair fixed tests exactly whether that same pair is achievable
under the threshold.  Rank-threshold feasibility is monotone, and the program
searches only thresholds where the edge set can change, so the returned
threshold is the smallest \(Q\) preserving the rank pair.  The fill-rate
tie-breakers are then optimized with the rank pair and threshold fixed, using
the exact rational comparison described below.  Each phase optimizes the next
component over the feasible set left by all previous phases.
\end{proof}

\section{Worst Rank Minimization for Fair Mode}

For a threshold \(T\), keep only edges satisfying
\[
  r_{ij} \leq T.
\]
Feasibility is monotone in \(T\).  The \texttt{fair} mode uses a lower-bound
max-flow feasibility test and binary search over the compressed rank-threshold
candidate list.  This same test also provides lower bounds and feasibility
checks for other exact objectives.  A laboratory edge has lower bound
\(\ell_j\) and upper bound \(c_j\).  The
network groups students by the allowed laboratory set for the tested
threshold, then uses the standard lower-bound transformation:
an edge \(u\to v\) with lower bound \(a\) and upper bound \(b\) is replaced by
capacity \(b-a\), and node balances are updated by
\[
  balance[u] \mathrel{-}= a,\qquad balance[v] \mathrel{+}= a.
\]
Then a super source and super sink are added, and feasibility is equivalent to
saturating all positive balance demands.

\begin{theorem}
The threshold returned by the search is the minimum possible value of \(Q(x)\).
\end{theorem}

\begin{proof}
For a fixed threshold \(T\), the transformed flow network is feasible if and
only if there exists an assignment satisfying all hard constraints using only
edges with \(r_{ij}\leq T\).  If a threshold is feasible, every larger
threshold is also feasible because it only adds edges.  Therefore the feasible
thresholds form an interval, and binary search returns the smallest feasible
threshold.
\end{proof}

\section{Lexicographic Min-Cost Flow}

After the optimal threshold \(T^*\) is known, only edges with
\(r_{ij}\leq T^*\) are used.  The program stores the ordinary min-cost-flow
part as a three-component
tuple
\[
  C=(c_0,c_1,c_2),
\]
ordered lexicographically.  Arithmetic is componentwise and checked for signed
integer overflow.  This avoids Big-M weights and avoids priority collapse from
scalarization.

The cost of assigning student \(i\) to laboratory \(j\) is
\[
  (0,\ r_{ij},\ r_{ij}^2).
\]
Laboratory slots encode required minimum counts.  If a slot is required, it
has first component \(-1\); otherwise its first component is 0.  Therefore
minimizing the first component maximizes the number of required slots used.
When all required slots can be used, the remaining components determine the
rank-optimal solution.
The final average fill-rate component is handled separately by the exact
least-common-denominator comparison in Section~\ref{sec:exact-average-fill}.  This
keeps the hot min-cost-flow path smaller while preserving the full
five-component objective.

\begin{theorem}
With \(T^*\) fixed, the lexicographic min-cost flow minimizes \(R_1\), and
among all assignments with minimum \(R_1\), minimizes \(R_2\).
\end{theorem}

\begin{proof}
All feasible solutions after the required-slot component is satisfied have the
same first component.  The second component of total cost is exactly
\(\sum r_{ij}x_{ij}=R_1(x)\).  Lexicographic minimization therefore minimizes
\(R_1\).  Among solutions with the same second component, the third component
is exactly \(\sum r_{ij}^2x_{ij}=R_2(x)\), so it is minimized next.
\end{proof}

\section{Exact Fill-Rate Tie-Breaking}

Different objective modes use the same exact fill-rate machinery in different
orders.  In \texttt{rubric}, \texttt{satisfaction}, and \texttt{guarded}, the
solver first maximizes \(U_{\mathrm{avg}}\) after the rank target and \(Q\)
have been fixed, then finds the largest \(U_{\min}\) threshold that preserves
that same average-fill value.  In \texttt{fair} and \texttt{balanced}, the
minimum-fill component is the next tie-breaker before the final average-fill
tie-breaker.

The minimum-fill subproblem is exact.  Every possible minimum fill-rate value
has the form
\[
  q=\frac{a}{c_j}
\]
for some positive-capacity laboratory \(j\) and integer
\[
  1 \leq a \leq \min(c_j,N).
\]
The program enumerates these rational candidates, sorts them by exact
cross-multiplication, removes duplicates, and binary-searches the best
candidate.

For a candidate \(q\), laboratory \(j\) must receive at least
\[
  \left\lceil q c_j \right\rceil
\]
students.  The program marks exactly that many slots as required.  It then
runs the same lexicographic min-cost flow.  The candidate is accepted if all
required slots are used and the rank components \(R_1\) and \(R_2\) remain
equal to the previously proven optima.

\begin{theorem}
When the selected objective asks for minimum fill before average fill, the
selected candidate maximizes \(U_{\min}\) without worsening the already fixed
rank and threshold components.
\end{theorem}

\begin{proof}
The candidate set contains every value that can be achieved by
\(n_j/c_j\).  For a fixed \(q\), the slot requirements are exactly equivalent
to \(n_j/c_j\geq q\) for all positive-capacity laboratories.  Feasibility is
monotone in \(q\): if \(q\) is feasible, every smaller candidate is feasible.
Binary search over the sorted candidate set therefore returns the greatest
feasible \(q\).  The acceptance test additionally requires the already fixed
rank components to remain unchanged, so the earlier optima are preserved.
\end{proof}

\section{Exact Average Fill-Rate Tie-Breaking}
\label{sec:exact-average-fill}

Whenever a mode needs to maximize \(U_{\mathrm{avg}}\) under previously fixed
components, the solver does so without using a floating-point or scaled
integer approximation.  Let
\[
  P=\{j\mid c_j>0\}
\]
and define the conceptual per-laboratory common denominator
\[
  L=\operatorname{lcm}\{c_j\mid j\in P\}.
\]
For each positive-capacity laboratory \(j\), define
\[
  W_j=\frac{L}{c_j}.
\]
Then
\[
  U_{\mathrm{avg}}(x)
  =\frac{1}{|P|}\sum_{j\in P}\frac{n_j}{c_j}
  =\frac{1}{|P|L}\sum_{j\in P}n_jW_j.
\]
Because \(|P|\) and \(L\) are fixed for the input instance, maximizing
\(U_{\mathrm{avg}}\) is exactly equivalent to maximizing the integer value
\[
  \sum_{j\in P} n_j W_j.
\]

For implementation, laboratories with the same positive capacity are grouped.
Let
\[
  C=\{d\mid \exists j,\ c_j=d,\ d>0\}
\]
be the set of distinct positive capacities, and let
\[
  B_d=\{j\mid c_j=d\}.
\]
Then
\[
  \sum_{j:c_j>0}\frac{n_j}{c_j}
  =
  \sum_{d\in C}\frac{\sum_{j\in B_d}n_j}{d}.
\]
Thus equal-denominator terms are combined before the least-common-denominator
integer comparison.  This transformation is algebraically exact and only
reduces the coefficient dimension.

The weights \(W_j\) can be much larger than 64-bit integers, especially when
there are many laboratories.  The program therefore stores them with a small
fixed-limb unsigned big-integer type.  After equal-capacity grouping, one
weight is stored per distinct positive capacity rather than per laboratory.
The final min-cost flow phase does not store a huge integer on every graph
edge.  Instead, each laboratory-to-sink edge carries only the capacity-term
index and a coefficient \(-1\).  During shortest-path comparisons, the solver
compares the accumulated coefficient vectors by evaluating the corresponding
big-integer weighted sums exactly.

There is also a scalar fast path for ordinary rank-first objectives.  Let
\[
  U_s(x)=\sum_{j\in P} n_j W_j
\]
be the exact scaled fill reward and let \(T(x)\) be the already lower-priority
rank-square or dissatisfaction-square component.  If \(L\), all \(W_j\), and
the required arithmetic bounds fit in signed integer costs, the solver chooses
\[
  M > \max_{x,y}|U_s(x)-U_s(y)|
\]
and replaces the final pair \((T,-U_s)\) by the single integer
\[
  T M - U_s.
\]
If \(T(x)<T(y)\), then \(T(y)-T(x)\ge 1\), so the \(M\) gap cannot be reversed
by any possible fill-reward difference.  If \(T(x)=T(y)\), minimizing
\(TM-U_s\) is exactly the same as maximizing \(U_s\).  Therefore this fast
path preserves the same lexicographic order.  If any safety check fails, the
solver uses the big-integer comparison path described above.

The implementation uses two exact-preserving refinements to make this safety
test less conservative.  First, under a fixed max-rank threshold \(q\), edges
with rank greater than \(q\) are not inserted into the flow graph.  Overflow
and big-\(M\) bounds therefore need only consider rank costs for active ranks
\(1,\ldots,q\).  Second, the fill-reward difference bound is computed from
active capacity-limited reward slots.  Each positive-capacity laboratory
\(j\) contributes at most
\[
  \min(g_j,a_j)
\]
slots of reward \(W_j\), where \(g_j\) is its graph capacity and \(a_j\) is
the number of students with an active edge to \(j\) under the same threshold.
No assignment can send more than \(g_j\) students to \(j\), and no assignment
can send more than those \(a_j\) reachable students.  Every complete
assignment therefore selects exactly \(N\) slots from this capped multiset, so
the sum of the largest \(N\) slot rewards is an upper bound on \(U_s\), and
the sum of the smallest \(N\) slot rewards is a lower bound.  Their
difference is a valid \(D\) for the proof above and is no larger than the
coarser \(N(\max_j W_j-\min_j W_j)\) bound.

\begin{theorem}[Exact fill tie-breakers]
After the rank-side prefix of a selected objective mode has been fixed, the
solver applies the selected fill tie-breaker exactly.  For average-first modes
such as \texttt{rubric}, \texttt{satisfaction}, and \texttt{guarded}, it first
maximizes \(U_{\mathrm{avg}}\), then maximizes \(U_{\min}\) subject to
preserving that optimal \(U_{\mathrm{avg}}\).  For minimum-first modes such as
\texttt{fair} and \texttt{balanced}, it first maximizes \(U_{\min}\), then
maximizes \(U_{\mathrm{avg}}\) subject to preserving that optimal \(U_{\min}\).
\end{theorem}

\begin{proof}
At this point the earlier rank-side constraints are represented as hard
conditions: the selected rank threshold and the already optimal additive rank
costs must be preserved.  In an average-first phase, each laboratory-to-sink
unit of flow for laboratory \(j\) contributes the exact reward represented by
\(W_j\).  Thus the fill component of a full flow is
\[
  -\sum_{j\in P} n_j W_j.
\]
Minimizing this value is exactly the same as maximizing
\(\sum_j n_j W_j\), and by the equality above, exactly the same as maximizing
\(U_{\mathrm{avg}}\).  The subsequent minimum-fill phase adds the minimum
count vector corresponding to a candidate threshold and searches for the
largest threshold that can preserve the already optimal average-fill value.

In a minimum-first phase the order is reversed.  The solver first searches
minimum-count threshold vectors to find the largest achievable
\(U_{\min}\).  With that vector fixed, the same exact average-fill comparison
maximizes \(U_{\mathrm{avg}}\).  In both cases, weighted sums are compared
with integer arithmetic, so no rounding can change the ordering.
\end{proof}

\subsection{Why Exact Average Comparison Is Needed Only in Fill Phases}

The average fill-rate comparison is not needed while the solver is still
deciding an earlier lexicographic component.  If two partial costs differ in
rank sum, rank-square sum, worst rank, or an already selected minimum-fill
threshold, \(U_{\mathrm{avg}}\) cannot change that lexicographic decision.
Therefore the expensive big-integer comparison is used only in phases where
the selected objective mode actually compares average fill among assignments
whose earlier components are already tied.

For average-first modes, the exact comparison chooses the member of the
rank-side optimum set with maximum \(U_{\mathrm{avg}}\), and a later phase
chooses maximum \(U_{\min}\) without losing that value.  For minimum-first
modes, the minimum-fill value is fixed first, and the exact comparison then
chooses maximum \(U_{\mathrm{avg}}\) inside that smaller optimum set.  Such
cases are possible with large capacities, which is why the submitted solver
uses exact big-integer comparison instead of claiming that a fixed scale is
always safe.

\subsection{Weighted Exact Objective Mode}

The implementation also provides an optional weighted exact mode.  It
minimizes
\[
  w_1R_1+w_2R_2+w_3Q-w_4U_{\mathrm{avg}}-w_5U_{\min}+w_6O+w_7C,
\]
where \(O\) is the number of assignments outside the submitted preference
list and \(C\) is the number of students moved away from a supplied baseline
assignment.  The weights are nonnegative integers.  The mode uses \(R_2\), not
the standard deviation itself, because \(R_2\) is additive over
student--laboratory edges.  The reported standard deviation remains exact for
the selected assignment.

The non-additive terms \(Q\) and \(U_{\min}\) are handled outside the flow
network as thresholds.  The solver searches a two-dimensional space whose
axes are max-rank thresholds \(q\) and minimum-fill lower-bound vectors
induced by candidate ratios \(u\).  For a box
\[
  q\in[q_{\mathrm{lo}},q_{\mathrm{hi}}],
  \qquad
  u\in[u_{\mathrm{lo}},u_{\mathrm{hi}}],
\]
it solves the relaxed corner \((q_{\mathrm{hi}},u_{\mathrm{lo}})\) to obtain
the additive lower bound \(A(q_{\mathrm{hi}},u_{\mathrm{lo}})\).  The whole
box is pruned only when
\[
  A(q_{\mathrm{hi}},u_{\mathrm{lo}})
  + w_3q_{\mathrm{lo}} - w_5u_{\mathrm{hi}}
\]
cannot beat the incumbent weighted score.  This branch-and-bound is exact:
the rank constraint is only relaxed, the minimum-fill constraint is only
relaxed, and the non-additive terms are bounded optimistically.
Indeed, for any assignment \(x\) in the box,
\[
  q_{\mathrm{lo}}\le Q(x)\le q_{\mathrm{hi}},\qquad
  u_{\mathrm{lo}}\le U_{\min}(x)\le u_{\mathrm{hi}}.
\]
The relaxed corner \((q_{\mathrm{hi}},u_{\mathrm{lo}})\) has a superset of the
box's feasible assignments, so
\[
  A(q_{\mathrm{hi}},u_{\mathrm{lo}})\le A(x).
\]
Because all weights are nonnegative,
\[
  w_3 q_{\mathrm{lo}}\le w_3 Q(x),
  \qquad
  -w_5 u_{\mathrm{hi}}\le -w_5 U_{\min}(x).
\]
Adding these inequalities gives a lower bound on the full weighted score of
every assignment in the box, so pruning cannot remove a better solution.
When the common denominator \(L\) and the scaled costs fit safely in signed
64-bit arithmetic, the \(U_{\mathrm{avg}}\) reward is folded into the ordinary
edge costs by minimizing
\[
  (|\{j:c_j>0\}|L)I - w_4\sum_{j:c_j>0} n_j\frac{L}{c_j},
\]
where \(I\) is the additive rank/outside cost.  If that scaling is not safe,
the solver falls back to the same least-common-denominator big-integer path
comparator as the default exact tie-breaker.  Thus \(U_{\mathrm{avg}}\) is part
of the primary weighted score and is never a floating-point approximation.
When hard targets bound the rank axis above by \(q_{\max}\), the scalar
safety check for this weighted average-fill path only needs to consider rank
costs up to \(q_{\max}\), because no weighted corner solve can insert edges
with larger ranks.

Average-fill hard targets are accepted only through exact bounded cases.  If
the existing laboratory lower bounds already imply
\[
  U_{\mathrm{avg}}\ge \alpha,
\]
then adding the target does not change the feasible region.  If average fill
is constant over complete assignments, for example when all positive
capacities are equal, the target either rejects the whole instance as
infeasible or again leaves the selected objective's feasible region unchanged.
For supported single objectives, the remaining nontrivial cases use a bounded
count-vector engine.  Let \(L\) be the common denominator of the positive
capacities and define
\[
  U_s(n)=\sum_{j:c_j>0} n_j\frac{L}{c_j}.
\]
The hard target \(U_{\mathrm{avg}}\ge\alpha\) becomes one integer inequality on
\(U_s(n)\).  The implementation enumerates laboratory count vectors satisfying
all lower bounds, graph-capacity upper bounds, \(\sum_j n_j=N\), and this
integer resource inequality.  For each retained vector, using \(n_j\) as the
minimum count for laboratory \(j\) fixes the counts exactly because the lower
bound sum is \(N\).  The solver then runs the appropriate exact min-cost-flow
subproblem for the selected rank-first or fair objective and compares all
count slices in that same objective order.  If the common-denominator scaling
or the exact count-vector bound is too large, the run is rejected rather than
silently using a heuristic.
Equivalent minimum-fill thresholds are compressed by their induced
lower-bound-count vectors.  Max-rank thresholds are compressed to the ranks
that can actually change the student--laboratory edge set, together with the
forced-placement sentinel rank \(M+1\).  The best candidate is selected by
exact integer comparison of the full weighted score.  This is an exact
optimizer for the stated weighted objective; it is not merely selecting among
a small portfolio of precomputed objectives.  The current implementation also
seeds the incumbent with an exact rank/outside/change-only weighted solution
evaluated under the full score, and processes search boxes in increasing
lower-bound order with a binary heap.  Both choices only improve pruning and
do not change the set of candidate assignments or the proof of optimality.
Relaxed-corner cache entries are allocated lazily, so unvisited
\((q,u)\) pairs cost only a pointer-table slot rather than a full
big-integer score and copied solution.

\section{Implementation Notes}

The implementation is a single C11 file, \texttt{assignment.c}.  The sequential
solver core uses portable C data structures and arithmetic.  Optional terminal
confirmation and process-parallel portfolio execution use POSIX facilities
when they are available; Windows/MinGW builds fall back to deterministic serial
portfolio solving.  Important implementation choices are:
\begin{itemize}[leftmargin=2em]
  \item Dinic's algorithm for lower-bound feasibility.
  \item Successive shortest augmenting paths with potentials for min-cost
        flow.
  \item Layered initial potentials on fresh assignment graphs.  Before any
        augmentation, positive residual capacity appears only on forward
        source--student/group, student/group--laboratory, and
        laboratory--sink edges, so topological relaxation gives the same
        initial distances as the general queue method.  If this layer
        condition is violated, the implementation falls back to the general
        initialization.
  \item A custom 4-ary indexed heap whose keys are lexicographic cost tuples.
  \item Exact rational comparison for fill-rate candidates using an
        overflow-safe two-word product.
  \item Exact average fill-rate comparison in the final tie-break phase using
        fixed-limb unsigned big integers, least-common-denominator integer
        weights, and equal-capacity bucket compression.
  \item Exact pre-reservation of adjacency-list capacities before constructing
        each min-cost flow graph.
  \item A lower-bound priority queue for weighted-exact branch-and-bound boxes,
        so the most promising exact search region is examined first.
  \item Lazy weighted-exact relaxed-corner cache allocation; this keeps the
        exact cache but avoids eagerly allocating one large entry for every
        possible threshold pair.
  \item UTF-8 laboratory names are treated as byte strings; no character-level
        decoding is required for equality.
  \item Laboratory names are indexed with an open-addressed hash table, so
        parsing full preference lists does not repeatedly scan all laboratories
        with \texttt{strcmp}.
  \item Large input files are read by a buffered \texttt{fread()} line reader
        that preserves line-level validation semantics.
  \item Student duplicate detection uses an open-addressed hash table, and
        per-student duplicate preferences use generation stamps.
  \item Allowed-set groups are used in the threshold feasibility search, and
        active rank-row groups are used in the general min-cost-flow graph and
        expanded back to student-wise output after optimization.
  \item A non-fatal \texttt{try\_solve\_with\_minimum\_counts} path lets
        threshold searches use the min-cost-flow result itself as the
        feasibility test, avoiding a redundant Dinic graph when exact
        optimality does not require a separate precheck.
  \item Optional profiling counters record Dinic calls, min-cost-flow calls,
        exact-average comparisons, group builds, threshold candidates, graph
        sizes, and phase CPU timings.
\end{itemize}

\subsection{Output Files}

The main output file keeps the simple assignment format.  The
first line contains the number of students \(N\).  Each following line
contains one student id and the assigned laboratory name:
\[
\begin{array}{l}
\texttt{student\_count}\\
\texttt{student\_id assigned\_laboratory\_name}\\
\texttt{student\_id assigned\_laboratory\_name}\\
\cdots
\end{array}
\]
Therefore the main output has exactly \(N+1\) lines.  This is the only file
written in the default submission mode.

If the optional \texttt{--reports} flag is specified, the program writes
sidecar files.  The first is named
\[
  \texttt{OUTPUT\_FILE.metrics.txt}.
\]
This file contains the numerical values of the five evaluation metrics:
average assigned rank, rank standard deviation, maximum assigned rank,
average fill rate, and minimum fill rate.  It also contains the rank sum,
rank-square sum, per-laboratory assigned counts, and CPU-time measurements
from the C standard \texttt{clock()} function.  The extra file does not change
the assignment-list format used by external evaluators.

The second sidecar file is named
\[
  \texttt{OUTPUT\_FILE.by\_lab.txt}.
\]
It lists each laboratory, its assigned count, capacity, fill rate, fill
percentage, and the student ids assigned there with their assigned ranks.  A
student assigned outside the submitted preference list is marked as
\texttt{outside}.

The remaining sidecars are
\[
  \texttt{OUTPUT\_FILE.by\_student.tsv}
\]
and
\[
  \texttt{OUTPUT\_FILE.outside\_preferences.tsv}.
\]
They provide a tab-separated student-wise view and a focused list of students
assigned outside their submitted preference list.  These reports support
faculty and student review while preserving the concrete student-wise output
format used by external evaluators.

The report mode also writes
\begin{center}
  {\small\path{OUTPUT_FILE.reasons.tsv}},
\end{center}
a compact per-student explanation label, including the first-choice demand,
capacity, and assigned count used to justify congestion explanations.

When \texttt{--explain-student} or \texttt{--try-lock} is specified, the
program also writes
\begin{center}
  {\small\path{OUTPUT_FILE.explain.tsv}}.
\end{center}
This file is an exact counterfactual explanation for one requested
student--laboratory lock.  The solver adds that lock, re-optimizes the
remaining assignment under the selected objective, and reports feasibility,
metric deltas, and the number of changed students.
It is intentionally a single-objective tool: the program rejects
counterfactual explanation combined with portfolio mode, because portfolio
mode chooses among several objective orders and the re-solve would otherwise
use a different decision rule from the selected assignment.

When \texttt{--profile} is specified, the program writes
\begin{center}
  {\small\path{OUTPUT_FILE.profile.tsv}},
\end{center}
with internal solver counters and phase timing measurements.  This file is not
part of the assignment output; it is a reproducibility and tuning aid used to
confirm whether an optimization removes actual work without changing the exact
objective.

For adjustment runs with a previous assignment, the program writes
\begin{center}
  {\small\path{OUTPUT_FILE.adjustment_delta.tsv}},
\end{center}
which compares baseline and adjusted metrics and lists moved students.
Portfolio mode solves a light exact set of rank-oriented objective modes, writes
\begin{center}
  {\small\path{OUTPUT_FILE.portfolio.tsv}},
\end{center}
and also writes candidate assignment files such as
\path{OUTPUT_FILE.rubric.txt}, \path{OUTPUT_FILE.satisfaction.txt}, and
\path{OUTPUT_FILE.fair.txt}.  A deep portfolio option also includes heavier
exact fill-focused, fill-convex, and weighted-exact candidates.  The metrics
and student reports include dissatisfaction summaries computed from the
selected rank-cost model, so the rank-oriented and satisfaction-oriented
candidates can be compared side by side.  The portfolio table includes the
normalized components of the local recommendation score: average-rank
component, rank-standard-deviation component, max-rank component,
average-fill deficit, and minimum-fill deficit.  This score is only a local
selection proxy among generated candidates; an external relative-rank score is
relative to all submitted programs.  The table also records per-row
\texttt{strengths} and \texttt{weaknesses} fields, which label the metrics
where that generated candidate is relatively best or worst among the generated
portfolio.  When \texttt{--portfolio-summary-only} is
specified, the program keeps the final assignment, the portfolio table, and
ordinary reports but removes per-candidate assignment sidecars after they have
been read.

The C program can also run portfolio candidates as independent child
processes with \texttt{--jobs}.  The repository includes a Python wrapper with
the same purpose.  This process-level parallelism does not alter the
mathematical objectives or the selected candidate metrics; it only overlaps
independent exact solver invocations on multi-core machines.  The C
implementation uses POSIX child processes when available and falls back to
serial solving on Windows/MinGW.  In parallel mode, reported solver CPU time
is the sum of candidate CPU times; the observed wall-clock wait time can be
smaller because the candidates run concurrently.

\subsection{Operational Constraints and Adjustment}

The optional constraints file supports hard constraints
\texttt{lock student lab}, \texttt{forbid student lab},
\texttt{allow student lab...}, and \texttt{capacity lab value}.  These are not
heuristics.  They restrict the feasible flow network, so the solver returns an
exact optimum among assignments satisfying the stated constraints.  With
\texttt{--base-assignment} and \texttt{--change-penalty}, moving a student away
from the supplied baseline adds an edge cost, which favors stable adjustments
without changing the default output format.

\subsection{Efficiency-Preserving Optimizations}

The code contains several optimizations that do not change the mathematical
objective.

First, laboratory-to-sink slots with the same cost are merged.  Earlier, a
laboratory with capacity \(c_j\) could be represented by \(c_j\) separate
unit-capacity edges.  This is unnecessary because all required slots have the
same tuple cost and all optional slots have the same tuple cost.  Therefore
the implementation uses at most two edges:
\[
  \text{required edge capacity} = m_j,\qquad
  \text{optional edge capacity} = c_j-m_j,
\]
where \(m_j\) is the current lower requirement for laboratory \(j\).  Sending
\(k\) units through a capacity edge costs exactly the same as sending those
\(k\) units through \(k\) identical unit edges, so the feasible assignments
and lexicographic optimum are unchanged.  The edge count is reduced
substantially in dense cases.

Second, Dijkstra's priority queue is an indexed heap.  Each vertex appears in
the heap at most once, and a better distance performs a decrease-key update
instead of inserting a duplicate record.  The shortest paths are identical to
ordinary Dijkstra with nonnegative reduced costs, but stale heap entries are
avoided.

Third, if the total graph capacity is exactly \(N\), then every feasible
solution has fixed laboratory counts.  The solver can skip the fill-rate
search because \(U_{\min}\) and \(U_{\mathrm{avg}}\) are already determined by
the hard constraints.  Also, when the accepted \(U_{\min}\) lower bounds sum
to \(N\), or when all positive capacities are equal, the average fill-rate
tie-breaker cannot change the objective; the feasible solution already found
during the \(U_{\min}\) search is reused.

Fourth, the hot min-cost-flow edge cost stores only the three integer
components that are used before the exact average-fill phase.  Earlier
versions kept a fourth scaled tie-break field in the same struct.  After
switching \(U_{\mathrm{avg}}\) to exact least-common-denominator comparison, that
field became redundant.  Removing it reduces edge memory traffic and speeds
every min-cost-flow run without changing the mathematical objective.

Fifth, in the exact average-fill phase, laboratories with the same positive
capacity are represented by one capacity term.  This reduces big-integer
weight construction, coefficient-vector copying, and exact heap-key
comparison from the number of positive-capacity laboratories to the number of
distinct positive capacities.  The transformation is exactly the
equal-denominator identity proved above.

Sixth, the threshold feasibility search groups students by allowed laboratory
set for the tested threshold.  The general min-cost-flow path then groups
students with identical active rank rows for the fixed worst-rank bound.  A
group node has capacity equal to the number of students in the group, and its
edges to laboratories use the representative student's rank costs.  The
grouped flow is expanded back to concrete student ids after optimization.
This reduces the graph from \(N\) student nodes to \(G\leq N\) group nodes
without changing any objective value.

Seventh, repeated ungrouped min-cost-flow solves reuse an active-arc template
when the problem pointer, constraint pointer, and rank threshold are unchanged.
The template contains only the student--laboratory arc list and adjacency
reservation counts.  It does not contain edge costs, residual capacities,
reverse edges, or laboratory lower-bound sink edges; all of those are rebuilt
for each solve.  Therefore the reuse cannot transfer residual flow state or an
objective-specific cost from one threshold check to another.  It only removes a
duplicate scan of the active edge set.

Eighth, the rubric objective uses monotonicity in the minimum-fill threshold
after the optimal average fill-rate has been fixed.  Let \(A(u)\) be the
maximum average fill-rate achievable while preserving the already fixed
rank-cost target and requiring \(U_{\min}\ge u\).  If \(u_1\le u_2\), the
feasible set for \(u_2\) is a subset of the feasible set for \(u_1\), so
\(A(u_1)\ge A(u_2)\).  The implementation first computes \(A(0)\), then finds
the largest lower-bound vector that still achieves that same value using a
short sequential prefix followed by exponential and binary search.  Every
accepted point is still solved by the exact flow optimizer, so this only
removes redundant threshold probes.

Ninth, if all students have identical rank rows, the program uses a closed
count-based solver.  The check costs only \(O(NM)\) comparisons, which is the
same order as reading the rank table once.  When the check succeeds, each
student is interchangeable with every other student, so the assignment
objective depends only on laboratory counts.  The solver first fixes the
rank-optimal count totals by rank group.  Lower-rank groups are filled before
higher-rank groups because each additional unit in a lower-rank group reduces
\(R_1\), and if \(R_1\) ties, cannot increase \(R_2\).  Inside the one possible
frontier group with equal rank, the program performs the same exact rational
\(U_{\min}\) search on counts, then assigns any remaining units to the
smallest capacities first.  This last step exactly maximizes
\(\sum_j n_j/c_j\), because each additional unit in laboratory \(j\) has
marginal value \(1/c_j\).  For minimum-first objectives, average fill rate is
used only after the minimum fill rate has been fixed.  For rubric-style
average-first objectives, the implementation instead preserves the best
average-fill value while searching the minimum-fill threshold, as described
above.  Finally, the count vector is
expanded back to the student-wise output format.

Tenth, the ordinary min-cost-flow Dijkstra queue uses a radix heap only when
the current residual graph admits an exact unsigned scalar key.  Before each
augmentation, the solver scans positive-residual edges and computes their
reduced costs.  The radix path is enabled only if every reduced component
\((c_1,c_2,c_3)\) is componentwise nonnegative and if checked scales \(B_2\)
and \(B_1\) fit such that
\[
  K(c)=c_1B_1+c_2B_2+c_3
\]
preserves lexicographic order for every simple path in the current residual
graph.  A one-unit difference in \(c_2\) is larger than the largest possible
\(c_3\) path difference, and a one-unit difference in \(c_1\) is larger than
the largest possible lower-component path difference.  Thus Dijkstra's
settled distances have monotone unsigned integer keys, satisfying the radix
heap invariant.  If any sign or overflow check fails, the previous binary heap
is used without changing the path costs.

\section{Complexity}

Let \(E\) be the number of forward edges in the final min-cost flow network.
In the worst case, student-laboratory edges are dense:
\[
  E = O(NM + N + M).
\]
This bound uses the compressed laboratory-to-sink representation above.  With
\(N\leq1024\) and \(M\leq256\), this remains practical in C, and larger
instances remain practical when the effective edge count is controlled by
active grouping or repeated preferences.  The min-cost flow sends at most
\(N\) units of flow.  Each augmentation performs an
indexed-heap shortest path computation.  The exact \(U_{\min}\) phase adds a
binary search over at most
\[
  \sum_j \min(c_j,N)
\]
rational candidates.  Under the target benchmark limits, this is at most
\(MN\le 256\cdot1024=262144\); for larger OSS inputs, the same polynomial
bound scales with the supplied \(M\) and \(N\).
The exact \(U_{\mathrm{avg}}\) phase is run only once after the previous
objectives are fixed, and is skipped when the average fill rate is already
fixed by the counts or by uniform positive capacities.  Let \(P\leq M\) be
the number of positive-capacity laboratories, let \(K\leq P\) be the number
of distinct positive capacities, and let \(B\) be the number of base-\(2^{31}\)
limbs needed for the least-common-denominator weights \(W_d\).  Exact
comparison of the average-fill component costs \(O(KB)\) only when the first
three integer cost components tie.  This keeps the exact rational tie-breaker
local to the final min-cost flow instead of slowing all previous feasibility
and \(U_{\min}\) searches.

If \(G_Q\) is the number of distinct active rank rows among students for the
fixed worst-rank bound \(Q\), the grouped min-cost-flow graph has \(G_Q\)
group nodes and \(O(G_QM+G_Q+M)\) forward edges in the dense case.  In the
threshold search, if \(G_T\) is the number of distinct allowed laboratory sets
for a tested threshold \(T\), the feasibility graph has \(G_T\) group nodes
instead of \(N\) student nodes.  Since both \(G_Q\) and \(G_T\) are at most
\(N\), grouping cannot increase the graph size and can substantially reduce it
when many students submit similar preference orders.

For the identical-rank fast path, the flow network is not built.  The dominant
work is sorting at most \(M\) laboratories and sorting at most
\(\sum_j \min(c_j,N)\) rational fill-rate candidates:
\[
  O(NM + M\log M + NM\log(NM)).
\]
Under repeated-rank inputs this is much smaller than repeated shortest-path
min-cost flow and is typically dominated by input and output time.

The implementation stores counts and graph indices in standard C \texttt{int}
values, while large arrays use dynamically allocated storage.  Therefore
larger-than-benchmark inputs are not rejected merely for exceeding
\(256\) laboratories or \(1024\) students, but the effective edge count still
grows with the number of active groups times the number of laboratories.
That product, together with available memory, is the main scaling limit.

\section{Verification}

The following checks were executed:
\begin{itemize}[leftmargin=2em]
  \item \texttt{gcc -std=c11 -O2 -Wall -Wextra -Wshadow -Wconversion -pedantic -Werror}
  \item \texttt{python3 -m pytest -q}
  \item AddressSanitizer and UndefinedBehaviorSanitizer test run in which the
        pytest build hook is forced to compile the sanitizer-instrumented
        binary.
  \item Small-instance exhaustive search tests.  These enumerate every
        assignment and verify that the solver's selected objective matches the
        true optimum.  Weighted exact tests compare the exact scalar score
        itself, not an unrelated tie-break tuple.
  \item Maximum-scale stress input with \(N=1024\) and \(M=256\).
\end{itemize}

\subsection{Dataset Suite}

The repository contains the following concrete datasets under
\texttt{datasets/}.  Each dataset has a \texttt{\_labs.txt} file and a
\texttt{\_prefs.txt} file.

\begin{center}
\begin{tabular}{lll}
\toprule
Dataset & Purpose & Optimality evidence \\
\midrule
\texttt{synthetic\_sample} & Handout-format sample & General flow theorem \\
\texttt{forced\_outside} & Forced outside preference & Exhaustive search \\
\texttt{zero\_capacity} & Zero-capacity laboratory & Exhaustive search \\
\texttt{rank\_variance} & Rank-sum and variance tie & Exhaustive search \\
\texttt{fill\_rate} & Minimum fill-rate tie & Exhaustive search \\
\bottomrule
\end{tabular}
\end{center}

For every small dataset, the test suite enumerates all possible assignments,
filters exactly the feasible assignments, and compares the solver result
against the exhaustive optimum for the selected objective mode.  This includes
the default rubric tuple
\[
  (R_1,\ R_2,\ Q,\ -U_{\mathrm{avg}},\ -U_{\min}),
\]
the fair tuple
\[
  (Q,\ R_1,\ R_2,\ -U_{\min},\ -U_{\mathrm{avg}}),
\]
satisfaction-cost tuples, and weighted-exact scalar scores.  This proves
optimality for those concrete datasets independently of the flow
implementation.  For larger datasets, exhaustive search is not computationally
possible; optimality follows from the mode-specific theorems above, which
apply to every valid input instance.  The test suite also contains a
huge-capacity case where a scaled integer average-fill tie-breaker would lose
the exact ordering; the solver still matches the exhaustive optimum because
the final comparison uses least-common-denominator big integers.

\subsection{Comparison with Exhaustive Search}

Exhaustive search is useful as a verifier for small inputs, but it is not a
candidate for the submitted algorithm.  It checks \(M^N\) assignments before
capacity pruning.  The following local timings compare a Python exhaustive
checker with the C solver on the same small \(M=4\) family:

\begin{center}
\begin{tabular}{rrrr}
\toprule
\(N\) & Assignments checked & Exhaustive search & C solver \\
\midrule
6  & 4,096     & 0.0092s & 0.0031s \\
8  & 65,536    & 0.2064s & 0.0027s \\
10 & 1,048,576 & 4.0502s & 0.0033s \\
\bottomrule
\end{tabular}
\end{center}

At the benchmark size, direct enumeration would require
\[
  256^{1024}=2^{8192}\approx 10^{2466}
\]
candidate assignments before pruning.  The flow formulation avoids this
enumeration by using the assignment structure directly.

\subsection{Performance}

The maximum-scale stress runs used \(N=1024\) and \(M=256\).  The table below
is generated by \texttt{scripts/run\_benchmarks.py}.  Wall time is measured by
the driver process, and CPU phase times are read from \texttt{--profile}.  The
values are machine-specific evidence rather than portable constants.

\begin{center}
{\small
\begin{tabular}{p{0.39\linewidth}rrp{0.24\linewidth}}
\toprule
Case & Wall time & Solver CPU time & Notes \\
\midrule
Uniform 256 labs / 1024 students, 8 preferences, rubric
  & 2.022s & 1.868s & reports/profile on \\
Popular-skewed 256 labs / 1024 students, 8 preferences, rubric
  & 2.456s & 2.400s & reports/profile on \\
Popular-skewed 256 labs / 1024 students, satisfaction
  & 2.553s & 2.456s & student-friendly rank costs \\
Popular-skewed 256 labs / 1024 students, fair
  & 0.896s & 0.850s & worst-rank first \\
Dense full preferences, 128 labs / 512 students, rubric
  & 0.609s & 0.598s & reports/profile on \\
Dense full preferences, 128 labs / 512 students, weighted evaluation
  & 0.738s & 0.726s
  & seed skipped on wide rank axis; 20 corner misses \\
Long-name full preferences, 256 labs / 1024 students, rubric
  & 2.712s & 2.458s & 65.9MB input; read CPU 0.227s \\
Popular-skewed 256 labs / 1024 students, weighted rank-only
  & 1.823s & 1.805s
  & integer-only fast path \\
Popular-skewed 256 labs / 1024 students, weighted evaluation
  & 8.288s & 8.133s
  & priority-queue 2D branch-and-bound; 9 corner-cache hits \\
Popular-skewed 128 labs / 512 students, exact counterfactual
  & 1.009s & 0.507s
  & selected solver CPU; counterfactual CPU 0.492s \\
\bottomrule
\end{tabular}
}
\end{center}

These cases are deliberately different: they include uniform preferences,
popular-skewed preferences, dense full-preference input, a 65.9MB long-name
input, weighted-exact threshold search, portfolio process parallelism, and
counterfactual re-optimization.  The
optimizations are therefore not just a special-case shortcut.  The slot
compression and indexed heap affect every min-cost flow solve; the capacity
and fill-rate shortcuts are used only when their mathematical preconditions
make the skipped work redundant.

On a 128-laboratory / 512-student popular-skewed case, the built-in light
portfolio took 1.895 seconds serially and 0.560 seconds with
\texttt{--jobs 4}.  With \texttt{--portfolio-summary-only}, the same light
parallel portfolio took 0.565 seconds and removed per-candidate sidecars after
reading them.  The deep portfolio took 3.837 seconds serially and 1.448 seconds
with \texttt{--jobs 4}.  This process parallelism is an
operational speedup only; each candidate is still produced by a normal exact
\texttt{assign\_labs} invocation.  The reported solver CPU time for parallel
portfolio mode is the sum of child-process solver CPU times, so it can be
larger than wall time.

\section{Conclusion}

The solver does not enumerate assignments for the documented objective modes.
Instead, for the stated assignment model and the selected structured
objective, it uses total unimodularity and integrality of network flow,
together with polynomial threshold enumeration and exact rational comparison,
to obtain a global optimum exactly.  This is not a claim of a universal
polynomial-time optimizer for arbitrary assignment objectives.  With only
black-box access to an arbitrary objective, exact optimization can require
exponential search: if some feasible assignment has not been queried, two
objectives can agree on every queried assignment while making that unqueried
assignment uniquely optimal in one case and not optimal in the other.

With a succinct unrestricted objective, the model can encode SAT even while
respecting the positive-capacity minimum-occupancy rule.  For each Boolean
variable \(x_i\), create two laboratories \(T_i\) and \(F_i\), each with
capacity 2.  Create one variable student \(X_i\) and two dummy students
\(D_{i,T}\) and \(D_{i,F}\).  Define the objective \(F_\varphi\) to score zero
exactly when \(D_{i,T}\) is assigned to \(T_i\), \(D_{i,F}\) is assigned to
\(F_i\), \(X_i\) is assigned to one of \(T_i,F_i\), and the truth assignment
induced by \(X_i\mapsto T_i\) as true and \(X_i\mapsto F_i\) as false satisfies
the Boolean formula \(\varphi\).  All other assignments score one.  The dummy
students satisfy the one-student lower bound in both truth-value laboratories,
and capacity 2 leaves room for the variable student.  Hence
\(\varphi\) is satisfiable if and only if \(\min_a F_\varphi(a)=0\).  A
polynomial-time exact optimizer for every such succinct objective would
therefore decide SAT in polynomial time, implying \(P=NP\).  The implemented
solver deliberately avoids this claim by supporting only structured objective
families with flow, threshold, rational-comparison, or separable-convex
structure.

The default mode is an exact Pareto-lexicographic flow solver for the
objective order:
\[
  R_1 \rightarrow R_2 \rightarrow Q \rightarrow U_{\mathrm{avg}}
  \rightarrow U_{\min}.
\]
Additional objective modes, hard constraints, baseline-change penalties, and
portfolio reports reuse the same exact flow core while preserving the required
student-wise output format.

\end{document}
